\documentclass[10pt,a4paper]{amsart}
\usepackage[margin=19mm]{geometry}
\usepackage{amsmath,amssymb,mathtools}
\usepackage[scr=rsfs,cal=euler]{mathalfa}
\usepackage{xfrac}
\usepackage[unicode,hidelinks,bookmarks=false]{hyperref}
\newcommand{\ie}{i.e.\ }
\newcommand{\cf}{cf.\ }
\newcommand{\resp}{respectively\ }
\newcommand{\Subsection}{\subsection}
\providecommand{\nobreakdash}{\nobreak\hbox{-}}
\setlength{\emergencystretch}{2em}
\multlinegap0pt
\usepackage{bm}
\usepackage{bbm}
\usepackage{dsfont}
\usepackage{xspace}
\usepackage{cancel}
\usepackage{mathtools}
\usepackage{amsthm}
\makeatletter
\@ifundefined{thm}{\newtheorem{thm}{Thm}}{}
\@ifundefined{lem}{\newtheorem{lem}[thm]{Lemma}}{}
\makeatother
\title[Double shuffle Lie algebra and special derivations]{Double shuffle Lie algebra and special derivations}
\author{Benjamin Enriquez, Hidekazu Furusho}
\date{Source: arXiv:2505.02265v2 (CC BY 4.0)}
\begin{document}
\maketitle

\noindent\emph{Source and licence.} Double shuffle Lie algebra and special derivations. Version arXiv:2505.02265v2 (CC BY 4.0), distributed under the Creative Commons Attribution 4.0 International licence, as recorded in the arXiv licence metadata for this exact version and shown on the abstract page https://arxiv.org/abs/2505.02265v2. Full rights notice: licenses/arxiv-2505.02265-CC-BY-4.0.txt.\par\smallskip

\noindent\textit{Editorial scope.} Counted unit: the Lem whose source label is \texttt{lem:exact:seq:0505} in the article (source0.tex lines 11277--11284, source label \texttt{lem:exact:seq:0505}), delivered together with the article's own complete original proof of it (source0.tex lines 11286--11443, one \texttt{proof} environment of 158 lines). No mathematics is written, repaired or re-derived: statement and proof are the article's own text line for line. Required context retained uncounted: lem:comm:alg (lines 11262--11270). Every definition, display and label of those lines is retained unchanged. Omitted from this package and not redistributed: 1--11261, 11271--11276, 11285--11285, 11444--14765. Nothing inside a retained excerpt is omitted or altered: every retained body is delivered exactly as the article's lines are written, so no omission receipt of any kind accompanies this package. Citations: the retained excerpts carry no citation command at all, so no bibliography entry of the article is transcribed and no reference is redistributed. Reference adaptation: none was needed and none was made; every reference inside the delivered unit resolves inside it, so no unresolved reference or citation is produced. Printed slips kept literal: no printed word, symbol, hypothesis or number of the counted statement or of its proof was altered. Layout and class: the article's own class is \texttt{amsart} with packages including amscd, epsfig, caption, mathrsfs, pgf, rotating, stmaryrd, tikz, MnSymbol, calc, xy, hyperref, stackengine, scalerel; the wrapper is the lane's pinned \texttt{amsart} editorial wrapper at 10pt on A4, which restates the theorem-like environments the retained text needs and adds the article's own macro definitions that the retained text needs (none), with extra wrapper packages (bm, bbm, dsfont, xspace, cancel, mathtools). The delivered numbering is the unit's own. Assets: no retained span contains a figure environment, a graphics-inclusion command, a PDF-page inclusion, a table or a TikZ picture; no image file is copied into this package, the package declares no asset dependency and no image file is redistributed with it. Licence: version arXiv:2505.02265v2 of this article is distributed under the Creative Commons Attribution 4.0 International licence, as recorded in the arXiv licence metadata for this exact version and shown on the abstract page https://arxiv.org/abs/2505.02265v2. A full rights notice is delivered at licenses/arxiv-2505.02265-CC-BY-4.0.txt. Characters: every retained excerpt is pure ASCII, so the delivered engine, XeLaTeX with its default Latin Modern text font, reports no missing character.
\begin{lem}\label{lem:comm:alg}
Let $n\geq1$ and $(a_1,\ldots,a_n),(b_1,\ldots,b_n)$ be non-colinear vectors in $\mathbb Q^n$. 
Let $x_1,\ldots,x_n$ be free commutative variables, $a:=\sum_i a_ix_i$, $b:=\sum_i b_ix_i$, so 
$a,b\in \mathbf k[x_1,\ldots,x_n]$. Then the sequence 
$$
\mathbf k[x_1,\ldots,x_n]\to\mathbf k[x_1,\ldots,x_n]^{\oplus2}\to\mathbf k[x_1,\ldots,x_n], 
$$
where the first map is $P\mapsto (a\cdot P,b\cdot P)$ and the second map is $(A,B)\mapsto b\cdot A-a\cdot B$, is exact. 
\end{lem}

\begin{lem}\label{lem:exact:seq:0505}
The sequence of $\mathbf k$-module morphisms 
\begin{equation}\label{exact:seq:0705}
\mathbf k[[e_0,f_\infty]]\oplus \hat V\to \hat V^{\oplus 2}\to \hat V, 
\end{equation}
where the first map is $(\varphi,\gamma)\mapsto (\varphi,\varphi)+(\gamma(e_0+f_\infty),(e_0+f_\infty)\gamma)$
and the second map is $(u,v)\mapsto (e_0+f_\infty)u-v(e_0+f_\infty)$, is exact. 
\end{lem}

\begin{proof}
For $\Sigma\subset\mathcal V$ a subset, let ${\mathtt M}(\Sigma)$ be the submonoid of $\mathcal V$ generated by $\Sigma$ 
(we denote the unit by $\emptyset$). Then ${\mathtt M}(e_0,e_1)\simeq\{e_0,e_1\}^*$ is a $\mathbf k$-module basis of $\mathcal V$. 
Let ${}_{1}{\mathtt M}^0_{1}$ be its submonoid consisting 
of the unit $\emptyset$ and of the elements starting with and ending in $e_1$; then there are set inclusions 
 ${}_{1}{\mathtt M}^0_{1}\subset{\mathtt M}(e_0,e_1)\subset\mathcal V$. 

There is a bijection ${\mathtt M}(e_0)\sqcup({\mathtt M}(e_0)\times({}_{1}{\mathtt M}^0_{1}\smallsetminus\{\emptyset\})
\times{\mathtt M}(e_0))\to {\mathtt M}(e_0,e_1)$, where ${\mathtt M}(e_0)\to {\mathtt M}(e_0,e_1)$ is the canonical
injection and ${\mathtt M}(e_0)\times({}_{1}{\mathtt M}^0_{1}\smallsetminus\{\emptyset\})
\times{\mathtt M}(e_0)\to {\mathtt M}(e_0,e_1)$ is induced by the product. This implies that  
the $\mathbf k[e_0]$-bimodule $\mathcal{V}$ is decomposed as the direct sum
$$
\mathcal V=\oplus_{w\in{} {}_{1}{\mathtt M}^0_{1}} \mathcal V_0(w),\quad \mathrm{where}\quad 
\mathcal V_0(w):=\mathrm{im}(
\textbf{k}%\mathbf k
[e_0]^{\otimes 2}\to \mathcal V,a\otimes b\mapsto a\cdot w\cdot b)
$$
is the $\mathbf k[e_0]$-subbimodule of $\mathcal V$ generated by $w$.  
When $w\neq\emptyset$ (resp. $w=\emptyset$), the map $\mathbf k[e_0]^{\otimes 2}\to \mathcal{V}_0(w)$ induced by 
$a\otimes b\mapsto a\cdot w\cdot b$ (resp. the map $\mathbf k[e_0]\to \mathcal V_0(\emptyset)$ induced by the inclusion 
$\mathbf k[e_0]\subset\mathcal V$) induces an isomorphism of $\mathbf k[e_0]$-bimodules, the 
$\mathbf k[e_0]$-bimodule (i.e. $\mathbf k[e_0^{(l)},e_0^{(r)}]=\mathbf k[e_0]^{\otimes 2}$-module) structure on 
$\mathbf k[e_0^{(l)},e_0^{(r)}]=\mathbf k[e_0]^{\otimes 2}$ (resp. on $\mathbf k[e_0]$) being the regular one
(resp. induced by the product $\mathbf k[e_0]^{\otimes 2}\to\mathbf k[e_0]$, i.e. the morphism 
$\mathbf k[e_0^{(l)},e_0^{(r)}]\to \mathbf k[e_0]$, $e_0^{(l)},e_0^{(r)}\mapsto e_0$). 

Applying to the above situation the algebra automorphism $(e_1\mapsto e_1,e_0\leftrightarrow e_\infty)$ of $\mathcal V$
(recall $e_\infty:=-e_0-e_1$), one defines a submonoid ${\mathtt M}(e_\infty,e_1)$ of $\mathcal V$,
its submonoid ${}_{1}{\mathtt M}^\infty_{1}$, and the set inclusions ${}_{1}{\mathtt M}_{1}^\infty\subset{\mathtt M}(e_\infty,e_1)
\subset \mathcal{V}$; we then have the $\mathbf k[e_\infty]$-bimodule decomposition
$$
\mathcal V=\oplus_{w\in{} {}_{1}{\mathtt M}^\infty_{1}} \mathcal{V}_\infty(w),\quad \mathrm{where}\quad 
\mathcal{V}_\infty(w):=\mathrm{im}(
\textbf{k}%\mathbf k
[e_\infty]^{\otimes 2}\to \mathcal{V},a\otimes b\mapsto a\cdot w\cdot b), 
$$
and the $\mathbf k[e_\infty]$-bimodule isomorphisms $\mathbf k[e_\infty]^{\otimes2}\to\mathcal V_\infty(w)$, $a\otimes b\mapsto 
a\cdot w\cdot b$ for $w\neq\emptyset$, and $\mathbf k[e_\infty]\to\mathcal V_\infty(\emptyset)$ induced by the inclusion 
$\mathbf k[e_\infty]\subset\mathcal V$. 

Recall the notation $e_i:=e_i\otimes1$, $f_i:=1\otimes e_i$ in the tensor square $V=\mathcal V^{\otimes2}$ for
$i\in\{0,1,\infty\}$. The tensor product of the above bimodule decompositions gives rise to a
$\mathbf k[e_0,f_\infty]$-bimodule decomposition 
\begin{equation}\label{direct;sum:V:0705}
V=\underset{(w,z)\in{} {}_{1}{\mathtt M}^0_{1}\times {}_{1}{\mathtt M}^\infty_{1}
}{\oplus}{V({w,z})}, \quad \mathrm{where}\quad 
{V}(w,z):=\mathcal V_0(w)\otimes\mathcal V_\infty(z).
\end{equation}
Then $V(w,z)$ is the $\mathbf k [e_0,f_\infty]$-subbimodule of $V$ generated by $w\otimes z$. 
The above bimodule isomorphisms induce  $\mathbf k [e_0,f_\infty]$-bimodule (i.e. $\mathbf k [e_0^{(l)},f_\infty^{(l)},e_0^{(r)},f_\infty^{(r)}]=\mathbf k [e_0,f_\infty]^{\otimes2}$-module) isomorphisms
\begin{equation}\label{case:240508:1833}
    {V}(w,z)\simeq
\begin{cases}
 \mathbf k[e_0,f_\infty] & \text{ if } w=z=\emptyset, \\ 
  \mathbf k[e_0,f_\infty^{(l)},f_\infty^{(r)}] &  \text{ if }  w=\emptyset, z\neq \emptyset,\\ 
  \mathbf k[e_0^{(l)},e_0^{(r)},f_\infty] & \text{ if }  w\neq\emptyset, z=\emptyset, \\ 
  \mathbf k[e_0^{(l)},e_0^{(r)},f_\infty^{(l)},f_\infty^{(r)}] & \text{ if }  w\neq\emptyset, z\neq\emptyset. \\ 
\end{cases}
\end{equation}
the $\mathbf k [e_0^{(l)},f_\infty^{(l)},e_0^{(r)},f_\infty^{(r)}]$-module structures on the algebras in the right-hand sides being 
induced by the algebra morphisms from  $\mathbf k [e_0^{(l)},f_\infty^{(l)},e_0^{(r)},f_\infty^{(r)}]$ to these algebras 
taking $(e_0^{(x)},f_\infty^{(x)})$ ($x\in\{r,l\}$) to $(e_0,f_\infty)$ in the first case, $(e_0,f_\infty^{(x)})$ in the second 
case, $(e_0^{(x)},f_\infty)$ in the third case,  and the identity algebra morphism in the last case. 

If $M$ is a $\mathbf k[e_0,f_\infty]$-bimodule, let $\alpha_M : M\to M^{\oplus 2}$ and 
$\beta_M : M^{\oplus 2}\to M$ be the maps $m\mapsto (m\cdot(e_0+f_\infty),(e_0+f_\infty)\cdot m)$
and $(m,n)\mapsto (e_0+f_\infty)\cdot m-n\cdot(e_0+f_\infty)$. Let also 
$\mathrm{diag} : \mathbf k[e_0,f_\infty]\to\mathbf k[e_0,f_\infty]^{\oplus2}$, $\varphi\mapsto 
(\varphi,\varphi)$ be the diagonal map. Then $\mathrm{diag}$ as well as $\alpha_M,\beta_M$ are 
$\mathbf k[e_0,f_\infty]$-bimodule morphisms. Moreover, the assignments $M\mapsto\alpha_M,\beta_M$ are 
compatible with direct sums, so 
\begin{equation}\label{direct:sums}
\alpha_{M\oplus N}=\alpha_M\oplus\alpha_N, \quad \beta_{M\oplus N}=\beta_M\oplus\beta_N
\end{equation}
for $M,N$ two $\mathbf k[e_0,f_\infty]$-bimodules. 

The formulas from the statement of the lemma induce a sequence of $\mathbf k$-module morphisms
\begin{equation}\label{sequence:mod}
\mathbf k[e_0,f_\infty]\oplus V\to V^{\oplus 2}\to V,
\end{equation}
which coincides (upon passing from $\mathbf k[e_0,f_\infty]$-bimodule structures to $\mathbf k$-module structures) 
with the following sequence of $\mathbf k[e_0,f_\infty]$-bimodule morphisms
\begin{equation}\label{sequence:bimod}
\mathbf k[e_0,f_\infty]\oplus V\stackrel{(i^{\oplus2}\circ \mathrm{diag})\oplus\alpha_V}{\to} V^{\oplus 2}\stackrel{\beta_V}{\to} V,   
\end{equation}
where $i : \mathbf k[e_0,f_\infty]\to V$ is the canonical injection. 
 
It follows from \eqref{direct:sums}, from the direct sum decomposition \eqref{direct;sum:V:0705}, and from the fact that 
the image of $i$ is contained in $V(\emptyset,\emptyset)$,  
that  \eqref{sequence:bimod} decomposes as the direct sum over $(w,z)\in 
{}_{1}{\mathtt M}^0_{1}\times {}_{1}{\mathtt M}^\infty_{1}\setminus\{(\emptyset,\emptyset)\}$
of the sequence of $\mathbf k[e_0,f_\infty]$-bimodule morphisms  
\begin{equation}\label{sequence:z:w}
V(w,z)\stackrel{\alpha_{V(w,z)}}{\to}V(w,z)^{\oplus 2}
\stackrel{\beta_{V(w,z)}}{\to}V(w,z)    
\end{equation}
and of the sequence of $\mathbf k[e_0,f_\infty]$-bimodule morphisms
\begin{equation}\label{sequence:empty:empty}
V(\emptyset,\emptyset)\oplus\mathbf k[e_0,f_\infty]
\stackrel{\alpha_{V(\emptyset,\emptyset)}\oplus\mathrm{diag}}{\to}V(\emptyset,\emptyset)^{\oplus 2}
\stackrel{\beta_{V(\emptyset,\emptyset)}}{\to}V(\emptyset,\emptyset)    
\end{equation}
corresponding to $(w,z)=(\emptyset,\emptyset)$. 

Depending on the values of $(w,z)\neq(\emptyset,\emptyset)$, the isomorphism 
\eqref{case:240508:1833} sets up an isomorphism between the sequence of $\mathbf k[e_0,f_\infty]$-bimodule morphisms
\eqref{sequence:z:w} and the following sequences of $\mathbf k[e_0^{(l)},e_0^{(r)},f_\infty^{(l)},f_\infty^{(r)}]$-module morphisms: 

$\bullet$ if $w\neq\emptyset$ and $z\neq\emptyset$, the sequence 
$$
\mathbf k[e_0^{(l)},e_0^{(r)},f_\infty^{(l)},f_\infty^{(r)}]
\to \mathbf k[e_0^{(l)},e_0^{(r)},f_\infty^{(l)},f_\infty^{(r)}]^{\oplus 2}
\to \mathbf k[e_0^{(l)},e_0^{(r)},f_\infty^{(l)},f_\infty^{(r)}], 
$$
where the first map is $P\mapsto ((e_0^{(r)}+f_\infty^{(r)})\cdot P,(e_0^{(l)}+f_\infty^{(l)})\cdot P)$ 
and the second map is $(A,B)\mapsto (e_0^{(l)}+f_\infty^{(l)})\cdot A-(e_0^{(r)}+f_\infty^{(r)})\cdot B$. 
This is an exact sequence by Lem. \ref{lem:comm:alg}. 

$\bullet$ if $w=\emptyset$ and $z\neq\emptyset$, the sequence 
$$
\mathbf k[e_0,f_\infty^{(l)},f_\infty^{(r)}]
\to \mathbf k[e_0,f_\infty^{(l)},f_\infty^{(r)}]^{\oplus 2}
\to \mathbf k[e_0,f_\infty^{(l)},f_\infty^{(r)}], 
$$
where the first map is $P\mapsto ((e_0+f_\infty^{(r)})\cdot P,(e_0+f_\infty^{(l)})\cdot P)$ 
and the second map is $(A,B)\mapsto (e_0+f_\infty^{(l)})\cdot A-(e_0+f_\infty^{(r)})\cdot B$. 
This is an exact sequence by Lem. \ref{lem:comm:alg}. 

$\bullet$ if $w\neq\emptyset$ and $z=\emptyset$, the sequence 
$$
\mathbf k[e_0^{(l)},e_0^{(r)},f_\infty]
\to \mathbf k[e_0^{(l)},e_0^{(r)},f_\infty]^{\oplus 2}
\to \mathbf k[e_0^{(l)},e_0^{(r)},f_\infty], 
$$
where the first map is $P\mapsto ((e_0^{(r)}+f_\infty)\cdot P,(e_0^{(l)}+f_\infty)\cdot P)$ 
and the second map is $(A,B)\mapsto (e_0^{(l)}+f_\infty)\cdot A-(e_0^{(r)}+f_\infty)\cdot B$. 
This is an exact sequence by Lem. \ref{lem:comm:alg}. 

When $(w,z)=(\emptyset,\emptyset)$, the isomorphism \eqref{case:240508:1833} sets up an isomorphism between the sequence 
of $\mathbf k[e_0,f_\infty]$-bimodule morphisms \eqref{sequence:empty:empty} and the following sequences of 
$\mathbf k[e_0^{(l)},e_0^{(r)},f_\infty^{(l)},f_\infty^{(r)}]$-module morphisms: 
$$
\mathbf k[e_0,f_\infty]^{\oplus 2}
\to \mathbf k[e_0,f_\infty]^{\oplus 2}
\to \mathbf k[e_0,f_\infty], 
$$
where the first map is $(P,Q)\mapsto (P+(e_0+f_\infty)\cdot Q,P+(e_0+f_\infty)\cdot Q)$ 
and the second map is $(A,B)\mapsto (e_0+f_\infty)\cdot (A-B)$. 
This is an exact sequence since $e_0+f_\infty$ is not a zero divisor in $\mathbf k[e_0,f_\infty]$. 

It follows that the sequences of $\mathbf k[e_0,f_\infty]$-bimodule morphisms  
\eqref{sequence:z:w} and \eqref{sequence:empty:empty} are all exact. This implies the exactness of  
the sequence of $\mathbf k[e_0,f_\infty]$-bimodule morphisms \eqref{sequence:bimod}, and therefore of  
the sequence of $\mathbf k$-module morphisms \eqref{sequence:mod}. The latter sequence of $\mathbf k$-module morphisms
is graded, and \eqref{exact:seq:0705} is its graded completion. The exactness of \eqref{sequence:mod} then 
implies that of \eqref{exact:seq:0705}. 
\end{proof}

\end{document}
